\cotefile{157-2}
\scanpage{243}
\mspage{189}

il faut prendre $Z$ sans élément final. \deleted{Mais}

Mais prenons $Z$ ens \unsure{ordonné}, \deleted{\unsure{et}} \unsure{avec} un objet initial $\unsure{\varphi|\emptyset|0}$ mais sans objet final, \deleted{\unsure{mais}} \inserted{et} prenons \inserted{pour} $z_0$ un élément \uline{maximal} de $Z$, de sorte que $Z'=z_0\backslash Z$ est \unsure{isom.} à $\{z_0\}$. Il faut donner \deleted{\unsure{alors}} un exemple où l'inclusion de $X_{z_0}$ dans $X$ n'est pas $\in W$. On va prendre

\diagram{\begin{tikzcd}[ampersand replacement=\&,column sep=small,row sep=tiny]
\& \deleted{\unsure{z}}a\ \deleted{\unsure{z_0}}\\
Z=\varphi \arrow[ur]\arrow[dr] \&\\
\& \deleted{\unsure{z}}b\ \deleted{\unsure{z}}=z_0
\end{tikzcd}}

donc la donnée de $X$ \unsure{cofibrée|cofibré} \unsure{scindée|scindé} sur $Z$ équivaut à un diagramme dans $\mathrm{Cat}$

\diagram{\begin{tikzcd}[ampersand replacement=\&,column sep=small,row sep=tiny]
\& X_a\\
X_0\arrow[ur,"\alpha"]\arrow[dr,"\beta"]\&\\
\& X_b\quad .
\end{tikzcd}}

Je pense que sous des conditions assez fortes (p.ex. si $\alpha,\beta$ sont des immersions ouvertes, ou des immersions \unsure{fermées} -- à vérifier !) la catégorie $X$ est isomorphe dans $\mathrm{Hot}$ à la somme amalgamée dans $\mathrm{Cat}$
\[
X_a\amalg_{X_0}X_b
\]
et il n'y \unsure{aucune} raison (sauf si on suppose $\alpha\in W$, et encore \ldots) que $X_b\unsure{\mathbin{-}|\to}X_a\amalg_{X_0}X_b$ soit dans $W$. Prenons $X_0=X_b=e$, $\beta=\mathrm{id}_e$ donc $\alpha$ est donné par la donnée d'un élément $t$ de $\mathrm{Ob}(X_a=T)$. La \deleted{catégorie $X$ est} \inserted{\unsure{somme} amalgamée} est $T$, qui n'est pas \unsure{isom.} à $e$ dans $\unsure{\mathrm{Hot}_{W}|\mathrm{Hot}_{w}}$,
% archive sheet number: 242
